\setcounter{section}{3}
\setlength{\parindent}{0pt}  % niente rientro
\setlength{\parskip}{6pt}
\section{Criticità}
Nel corso dello sviluppo del progetto Rinova il gruppo ha dovuto affrontare diverse criticità tecniche ed organizzative. La prima delle quali la variazione improvvisa dell'organico del gruppo, nella fase iniziale infatti, immediatamente successiva alla definizione degli obbiettivi e dei macro requisiti il terzo membro del gruppo ha dovuto abbandonare il progetto a causa di un cambio di percorso di studi, questo ha necessariamente portato ad un adattamento del perimetro del progetto, in modo da adattarne il carico di lavoro.
La seconda criticità è stata dovuta alla marcata incompatibilità degli orari accademici e da periodi di indisponibilità alternata dei membri: \textbf{Enrico Sartori} ha avuto una ridotta capacità operativa nella prima parte del semestre per esigenze personali, mentre \textbf{Stefano Pezzetta} ha dovuto limitare il contributo durante l'avvio della sessione invernale d'esame.
Questi problemi sono stati superati grazie all'adozione di una metodologia di comunicazione ben strutturata per via telematica, inoltre, durante i periodi di "fermo" di un membro, l'altro poteva portare avanti attività indipendenti, scongiurando colli di bottiglia e permettendo il riallinamento successivamente.

L'analisi del contesto normativo riguardante le Comunità Energetiche Rinnovabili (CER) si è rivelata più onerosa del previsto, la burocrazia confusa e la gestione complessa dei rapporti fra produttori e fornitori di servizi ha contribuito a rallentare notevolmente la fase di studio di fattibilità del progetto.

Dal punto di vista implementativo, la difficoltà maggiore ha riguardato il reperimento e la normalizzazione dei dati provenienti dagli strumenti di misurazione.

chiedi quali sono state le difficoltà tecniche!